% bayesfilter-genut-score-variance-problem-and-repair-note-2026-07-31.tex
%
% Standalone proposition--proof note.  Self-contained preamble; no macro
% redefinitions (MathDevMCP parser compatibility).  Companion artifacts:
%   docs/plans/bayesfilter-genut-score-computation-audit-result-2026-07-30.md
%   docs/plans/bayesfilter-genut-score-estimator-options-mathematical-note-2026-07-31.md
\documentclass[11pt]{article}

\usepackage[margin=1.05in]{geometry}
\usepackage{amsmath}
\usepackage{amssymb}
\usepackage{amsthm}
\usepackage{booktabs}
\usepackage[colorlinks=true,linkcolor=blue,citecolor=blue]{hyperref}

\newtheorem{proposition}{Proposition}[section]
\newtheorem{lemma}[proposition]{Lemma}
\newtheorem{corollary}[proposition]{Corollary}
\theoremstyle{definition}
\newtheorem{definition}[proposition]{Definition}
\newtheorem{assumption}[proposition]{Assumption}
\theoremstyle{remark}
\newtheorem{remark}[proposition]{Remark}

\title{The GenUT-Route Score-Variance Problem and Proposed Repairs:\\
A Proposition--Proof Note}
\author{BayesFilter research notes}
\date{2026-07-31}

\begin{document}
\maketitle

\begin{abstract}
The staged GenUT candidate filter computes a manual recursive
forward-sensitivity score that was audited (2026-07-30) to be the exact total
derivative of the executed finite likelihood scalar on its fixed branch.
Empirically that score has very large particle-seed variance on the Austria
SIR target.  Part~I of this note isolates the mathematical mechanism: the
carried parameter tangent obeys a linear recursion whose stage operators have
gain bounds controlled only by ridge and damping floors
(Propositions~\ref{prop:gsv-tangent-recursion}--\ref{prop:gsv-gn-gain}), and
the replicated cubature residual design injects, at every reset, residuals
whose per-axis kurtosis equals the state dimension $d$ and whose pairwise
co-kurtosis is zero (Proposition~\ref{prop:gsv-cubature-design-moments}); at
$d=18$ this is a large designed-in shape defect recycled through the
recursion.  Part~II also defines a fourth, experimental repair: a frozen
low-rank Tucker subspace in which the \emph{complete} third central moment
and fourth cumulant are matched without materializing dense $d^3$ or $d^4$
tensors (Section~\ref{sec:gsv-projected-cumulants}).  This construction is a
project derivation and an \texttt{extension\_or\_invention}; Zhao and Cui's
squared-TT recursion motivates low-rank compression but does not contain this
moment correction.  The other three repairs are an exactly whitened fixed Gaussian
residual design with correct first two moments and asymptotically Gaussian
higher moments
(Propositions~\ref{prop:gsv-whitened-design-identities}--\ref{prop:gsv-delta-method}),
a surrogate-force Hamiltonian Monte Carlo scheme that remains exactly
invariant for the executed finite target under an \emph{arbitrary}
deterministic force field
(Propositions~\ref{prop:gsv-leapfrog-volume}--\ref{cor:gsv-surrogate-force}),
and Fisher-identity score estimators whose recursions contain no derivative
of any transport or reset map
(Propositions~\ref{prop:gsv-fisher-identity}--\ref{prop:gsv-backward-recursion}).
Exactness boundaries and non-claims are stated throughout: nothing here
asserts an exact-posterior score, unbiasedness of the finite scalar, or any
empirical superiority.
\end{abstract}

\tableofcontents

\section{Scope, executed program, and standing assumptions}
\label{sec:gsv-scope}

This note concerns the opt-in staged candidate implemented in
\texttt{bayesfilter/highdim/cubature\_genut\_filter.py},
\texttt{higher\_moment\_contract\_e.py}, and
\texttt{ledh\_contract\_e\_reset\_tf.py}, documented in chapter
\texttt{ch32c\_entropic\_ot\_sinkhorn.tex} (in particular its higher-moment
Contract~E section and its total-score proposition, referred to below as the
\emph{chapter score proposition}).  All mathematical objects are finite:
$N$ particles, state dimension $d$, parameter dimension $P$, horizon $T$.

\begin{definition}[Executed finite program]
\label{def:gsv-program}
Fix observations $y_{1:T}$, innovation rows $e_{t,i}$, an initial-noise
matrix, a residual design $\Xi\in\mathbb{R}^{N\times d}$, and all numerical
controls (Sinkhorn $\varepsilon$ and iteration counts, Contract-E ridge
$\lambda>0$, correction strength $\rho\ge 0$, damping $\delta>0$, correction
counts, pairwise controls).  Given the carried equal-weight cloud
$X_{t-1}^{+}\in\mathbb{R}^{N\times d}$, one time step computes
\begin{align}
 x_{t,i}^{-} &= f_\theta\!\left(x_{t-1,i}^{+},\,e_{t,i}\right),
 \label{eq:gsv-transition}\\
 \ell_{t,i} &= \log g_\theta\!\left(y_t \mid x_{t,i}^{-}\right),
 \qquad
 \widehat\ell_t = \log\!\Big(\tfrac1N\textstyle\sum_{i=1}^N
 e^{\ell_{t,i}}\Big),
 \label{eq:gsv-increment}\\
 \alpha_{t,i} &= \frac{e^{\ell_{t,i}}}{\sum_{j=1}^N e^{\ell_{t,j}}},
 \label{eq:gsv-weights}\\
 X_t^{+} &= R_\theta\!\left(X_t^{-},\,\alpha_t\right),
 \label{eq:gsv-reset}
\end{align}
where $R_\theta$ is the composed entropic-OT row quotient, Contract-E
Cholesky reset with design $\Xi$, and finite higher-moment correction.  The
finite scalar and its score are
\begin{equation}
\label{eq:gsv-scalar}
 \widehat L^N(\theta)=\sum_{t=1}^T \widehat\ell_t(\theta),
 \qquad
 \text{score} \;=\; \frac{d\,\widehat L^N}{d\theta}\in\mathbb{R}^{P}.
\end{equation}
\end{definition}

\begin{assumption}[Fixed differentiable branch]
\label{ass:gsv-branch}
Throughout, $\theta$ lies in an open set on which: every Cholesky argument in
the program is strictly positive definite; every Sinkhorn row mass exceeds its
floor; the cost-scale maximum branch is locally constant; and no validity gate
changes state.  On this set every stage map of
Definition~\ref{def:gsv-program} is continuously differentiable in all of its
inputs, and the manual recursive tangent equals the total derivative of
\eqref{eq:gsv-scalar} (this is the chapter score proposition; it was verified
against the implementation line-by-line in the 2026-07-30 audit and is taken
as given here).
\end{assumption}

\begin{definition}[The two derivative objects]
\label{def:gsv-objects}
Object~A is the \emph{same-scalar score} $d\widehat L^N/d\theta$ of
\eqref{eq:gsv-scalar}.  Object~B is the \emph{exact score}
$\nabla_\theta \log p_\theta(y_{1:T})$ of the underlying state-space model.
$\widehat L^N$ is a biased approximation of $\log p_\theta(y_{1:T})$ at
finite $N$ (the OT/moment reset is not an unbiasedness-preserving
resampling), so Object~A is a biased estimator of Object~B even in exact
arithmetic.  The empirical problem addressed in this note is the
particle-seed \emph{variance} of Object~A; its bias against Object~B is a
separate question for which Section~\ref{sec:gsv-fisher} provides reference
estimators.
\end{definition}

\section{Part I: the problem}
\label{sec:gsv-problem}

\subsection{The tangent recursion and its accumulation}

Write $v_t=\operatorname{vec}\!\big(dX_t^{+}/d\theta_p\big)\in\mathbb{R}^{Nd}$
for one parameter direction $\theta_p$, and let
$\Phi_t(\cdot,\theta):\mathbb{R}^{Nd}\to\mathbb{R}^{Nd}$ denote the composed
one-step map $X_{t-1}^{+}\mapsto X_t^{+}$ of
\eqref{eq:gsv-transition}--\eqref{eq:gsv-reset} at fixed $e_t,y_t$.

\begin{proposition}[Tangent recursion and its solution]
\label{prop:gsv-tangent-recursion}
Under Assumption~\ref{ass:gsv-branch}, with
$J_t = D_X\Phi_t\big(X_{t-1}^{+},\theta\big)\in\mathbb{R}^{Nd\times Nd}$ and
$b_t = \partial_\theta\Phi_t\big(X_{t-1}^{+},\theta\big)\in\mathbb{R}^{Nd}$
evaluated on the executed path,
\begin{equation}
\label{eq:gsv-recursion}
 v_t = J_t v_{t-1} + b_t ,
\end{equation}
and consequently
\begin{equation}
\label{eq:gsv-recursion-solution}
 v_T=\sum_{s=1}^{T}\Big(\prod_{u=s+1}^{T}J_u\Big) b_s
 \;+\;\Big(\prod_{u=1}^{T}J_u\Big) v_0 ,
\end{equation}
with the empty product equal to the identity and matrix products ordered
$J_T J_{T-1}\cdots$.
\end{proposition}

\begin{proof}
Equation~\eqref{eq:gsv-recursion} is the chain rule for
$X_t^{+}=\Phi_t(X_{t-1}^{+}(\theta),\theta)$ on the fixed branch.  For
\eqref{eq:gsv-recursion-solution}, induct on $T$: the case $T=0$ is trivial;
if \eqref{eq:gsv-recursion-solution} holds at $T-1$, then
$v_T=J_T v_{T-1}+b_T
=J_T\sum_{s\le T-1}(\prod_{u=s+1}^{T-1}J_u)b_s
+J_T(\prod_{u\le T-1}J_u)v_0+b_T$,
which regroups to \eqref{eq:gsv-recursion-solution} at $T$.
\end{proof}

\begin{proposition}[Increment tangent]
\label{prop:gsv-increment-tangent}
Let $\ell_i(\theta)=h\big(\theta,x_i(\theta)\big)$ be continuously
differentiable and
$\widehat\ell=\log\big(N^{-1}\sum_i e^{\ell_i}\big)$,
$\alpha_i=e^{\ell_i}/\sum_j e^{\ell_j}$.  Then
\begin{equation}
\label{eq:gsv-increment-tangent}
 \frac{d\widehat\ell}{d\theta_p}
 =\sum_{i=1}^N \alpha_i\,\frac{d\ell_i}{d\theta_p},
 \qquad
 \frac{d\ell_i}{d\theta_p}
 =\partial_{\theta_p} h\big(\theta,x_i\big)
 +\nabla_x h\big(\theta,x_i\big)^{\!\top}\frac{dx_i}{d\theta_p}.
\end{equation}
Hence the score \eqref{eq:gsv-scalar} is a sum over $t$ of linear functionals
of the carried tangents $v_{t-1}$ plus explicit-$\theta$ terms.
\end{proposition}

\begin{proof}
Differentiate: $d\widehat\ell
=\big(\sum_i e^{\ell_i}\big)^{-1}\sum_i e^{\ell_i}\,d\ell_i
=\sum_i\alpha_i\,d\ell_i$; the second identity is the chain rule.  The final
statement follows because $x_i=x_{t,i}^{-}$ depends on $\theta$ through
$X_{t-1}^{+}$ (a linear functional of $v_{t-1}$ at first order) and through
the explicit transition dependence.
\end{proof}

\begin{corollary}[Accumulation bound]
\label{cor:gsv-accumulation}
For any submultiplicative matrix norm,
\begin{equation}
\label{eq:gsv-accumulation}
 \lVert v_T\rVert
 \le \sum_{s=1}^{T}\Big(\prod_{u=s+1}^{T}\lVert J_u\rVert\Big)
 \lVert b_s\rVert
 +\Big(\prod_{u=1}^{T}\lVert J_u\rVert\Big)\lVert v_0\rVert .
\end{equation}
If the per-step gains satisfy $\lVert J_u\rVert\le e^{\gamma}$ with
$\gamma>0$ on a stretch of steps, the bound on early contributions grows
exponentially in the stretch length.  At $\gamma=0$ the geometric bound is
generally $O(T)$, not summable; a strictly negative bound prevents exponential
accumulation under bounded forcing.  A statement such as
$\operatorname{Var}(\mathrm{score})=O(T)$ additionally requires suitable
stationarity, moment, and mixing assumptions on the random stages and forcing
terms, and is not implied by a nonpositive exponent alone.
\end{corollary}

\begin{proof}
Apply the triangle inequality and submultiplicativity to
\eqref{eq:gsv-recursion-solution}; the growth statements are immediate from
the geometric factors.
\end{proof}

\begin{remark}[What is and is not established here]
\label{rem:gsv-lyapunov}
Corollary~\ref{cor:gsv-accumulation} is an upper bound; it identifies the
quantity that should be measured to characterize the regime: realized
finite-horizon directional growth
$g_t=\lVert J_t w_t\rVert/\lVert w_t\rVert$ of propagated probe tangents
$w_t$, with $\widehat\gamma = T^{-1}\sum_t\log g_t$.  This is not
automatically the asymptotic top Lyapunov exponent: one probe can miss the
dominant direction and nonnormal transients can dominate.  Use multiple fixed
probes (or QR/subspace iteration) and treat the result as explanatory evidence,
not a promotion criterion.  The probe has zero explicit-$\theta$ source, but
its extra tangent columns still incur a measurable memory and tracing cost.
\end{remark}

\subsection{Stage-gain inventory}

The next three propositions bound the gains of the stages that dominate
$\lVert J_t\rVert$ in the implemented route.  Each formula below is exactly
the differentiated map verified in the 2026-07-30 code audit.

\begin{proposition}[Weight-normalization gain]
\label{prop:gsv-weight-gain}
Let $\alpha_i=e^{\ell_i}/\sum_j e^{\ell_j}$ and let dots denote derivatives
in one parameter direction.  Then
$\dot\alpha_i=\alpha_i\big(\dot\ell_i-\sum_j\alpha_j\dot\ell_j\big)$ and
\begin{equation}
\label{eq:gsv-weight-gain}
 \sum_{i=1}^N\lvert\dot\alpha_i\rvert
 \;\le\; 2\max_{i}\lvert\dot\ell_i\rvert
 \;\le\; 2\max_i\Big(\lvert\partial_\theta \ell_i\rvert
 +\lVert\nabla_x\ell_i\rVert_2\,\lVert\dot x_i\rVert_2\Big).
\end{equation}
\end{proposition}

\begin{proof}
Differentiating $\log\alpha_i=\ell_i-\log\sum_j e^{\ell_j}$ gives
$\dot\alpha_i/\alpha_i=\dot\ell_i-\sum_j\alpha_j\dot\ell_j$.  Then
$\lvert\dot\alpha_i\rvert\le\alpha_i\big(\lvert\dot\ell_i\rvert
+\max_j\lvert\dot\ell_j\rvert\big)\le2\alpha_i\max_j\lvert\dot\ell_j\rvert$,
and summing over $i$ uses $\sum_i\alpha_i=1$.  The second inequality is
Cauchy--Schwarz applied to \eqref{eq:gsv-increment-tangent}.
Interpretation: particle tangents enter weight tangents with gain of order
$\max_i\lVert\nabla_x\ell_i\rVert_2$, which is large exactly when the
likelihood is informative.
\end{proof}

\begin{proposition}[Contract-E affine gain]
\label{prop:gsv-contract-e-gain}
Let $\widetilde P,P_w\in\mathbb{R}^{d\times d}$ be symmetric with
$\widetilde P+\lambda I$ and $P_w+\lambda I$ positive definite for a fixed
$\lambda>0$, and let
$L_E L_E^{\top}=\widetilde P+\lambda I$,
$L_w L_w^{\top}=P_w+\lambda I$ be their lower-triangular Cholesky factors;
positive definiteness of $L_EL_E^{\top}$ makes $L_E$ nonsingular, so
$L_E^{-1}$ exists and $A=L_w L_E^{-1}$ is well defined.  Then, for
differentiable paths of these factors,
\begin{equation}
\label{eq:gsv-affine-tangent}
 \dot A=\big(\dot L_w-A\dot L_E\big)L_E^{-1},
\end{equation}
\begin{equation}
\label{eq:gsv-affine-gain}
 \bigl\lVert L_E^{-1}\bigr\rVert_2
 =\lambda_{\min}\!\big(\widetilde P+\lambda I\big)^{-1/2}
 \le\lambda^{-1/2},
 \qquad
 \lVert\dot A\rVert_2
 \le\big(\lVert\dot L_w\rVert_2+\lVert A\rVert_2\lVert\dot L_E\rVert_2\big)
 \,\lambda^{-1/2}.
\end{equation}
\end{proposition}

\begin{proof}
Differentiate $A=L_wL_E^{-1}$ with
$d(L_E^{-1})=-L_E^{-1}\,\dot L_E\,L_E^{-1}$:
$\dot A=\dot L_wL_E^{-1}-L_wL_E^{-1}\dot L_EL_E^{-1}
=(\dot L_w-A\dot L_E)L_E^{-1}$, which is \eqref{eq:gsv-affine-tangent}.
For a nonsingular square matrix $L$, the matrices $L^{\top}L$ and
$L L^{\top}$ are similar ($L^{\top}L=L^{-1}(LL^{\top})L$), hence share
eigenvalues, so
$\sigma_{\min}(L_E)^2=\lambda_{\min}(L_E^{\top}L_E)
=\lambda_{\min}(L_EL_E^{\top})
=\lambda_{\min}(\widetilde P+\lambda I)\ge\lambda$.
Thus $\lVert L_E^{-1}\rVert_2=\sigma_{\min}(L_E)^{-1}\le\lambda^{-1/2}$,
and the bound on $\lVert\dot A\rVert_2$ follows from
\eqref{eq:gsv-affine-tangent} by submultiplicativity.  At the production
value $\lambda=10^{-5}$ the worst-case factor is
$\lambda^{-1/2}\approx316$ per reset.
\end{proof}

\begin{proposition}[Damped Gauss--Newton coefficient gain]
\label{prop:gsv-gn-gain}
Let $J_a\in\mathbb{R}^{2\times2}$, $\delta>0$,
$M_a=J_a^{\top}J_a+\delta I\succeq\delta I$, $r_a=J_a^{\top}d_a$, and
$c_a=\rho\,M_a^{-1}r_a$ (the implemented correction coefficients); since
$M_a\succeq\delta I$ with $\delta>0$, the matrix $M_a$ is symmetric positive
definite, hence nonsingular, so $M_a^{-1}$ exists.  Then
\begin{equation}
\label{eq:gsv-gn-tangent}
 \dot c_a=\rho\,M_a^{-1}\big(\dot r_a-\dot M_aM_a^{-1}r_a\big),
 \qquad
 \lVert\dot c_a\rVert_2
 \le\frac{\rho}{\delta}\Big(\lVert\dot r_a\rVert_2
 +\frac{1}{\delta}\lVert\dot M_a\rVert_2\lVert r_a\rVert_2\Big).
\end{equation}
\end{proposition}

\begin{proof}
Differentiate $M_ac_a=\rho r_a$:
$\dot M_ac_a+M_a\dot c_a=\rho\dot r_a$, so
$\dot c_a=M_a^{-1}(\rho\dot r_a-\dot M_ac_a)
=\rho M_a^{-1}(\dot r_a-\dot M_aM_a^{-1}r_a)$.
Since $M_a\succeq\delta I$ is symmetric positive definite,
$\lVert M_a^{-1}\rVert_2\le\delta^{-1}$; apply it twice.  At the production
value $\delta=10^{-5}$ the worst-case outer factor is $10^{5}$, attained
only when $J_a$ degenerates; clouds near the moment-feasibility boundary
(empirical kurtosis close to $1+\text{skewness}^2$) drive $J_a$ toward
singularity.
\end{proof}

\subsection{The designed-in shape defect of the replicated cubature design}

\begin{proposition}[Cubature residual-design moments]
\label{prop:gsv-cubature-design-moments}
Let $d\ge1$, $M\ge1$, $N=2dM$, and let
$\Xi_c\in\mathbb{R}^{N\times d}$ consist of the $2d$ rows
$\sqrt{d}\,(\pm e_a)$, $a=1,\dots,d$, each repeated $M$ times.  Then, with
$r$ indexing rows,
\begin{align}
 \frac1N\sum_{r}\Xi_{ra}&=0,
 &\frac1N\sum_{r}\Xi_{ra}\Xi_{rb}&=\delta_{ab},
 \label{eq:gsv-cubature-first-second}\\
 \frac1N\sum_{r}\Xi_{ra}^{3}&=0,
 &\frac1N\sum_{r}\Xi_{ra}^{4}&=d,
 \label{eq:gsv-cubature-fourth}\\
 \frac1N\sum_{r}\Xi_{ra}^{2}\Xi_{rb}&=0\quad(a\neq b),
 &\frac1N\sum_{r}\Xi_{ra}^{2}\Xi_{rb}^{2}&=0\quad(a\neq b).
 \label{eq:gsv-cubature-cross}
\end{align}
The Gaussian reference values for the two fourth-order statistics are $3$
and $1$ respectively.
\end{proposition}

\begin{proof}
Coordinate $a$ is nonzero on exactly the $2M$ rows $\pm\sqrt d\,e_a$, taking
values $\pm\sqrt d$ in equal numbers.  Odd powers cancel in pairs, giving the
zero mean, zero third moment, and zero $\Xi_{ra}^2\Xi_{rb}$ (the latter also
because every row has a single nonzero coordinate).  For the variance,
$\tfrac1N\sum_r\Xi_{ra}^2=\tfrac{2M\cdot d}{2dM}=1$; off-diagonal second
moments vanish because no row has two nonzero coordinates, giving
\eqref{eq:gsv-cubature-first-second}.  For the fourth moment,
$\tfrac1N\sum_r\Xi_{ra}^4=\tfrac{2M\cdot d^{2}}{2dM}=d$.  Finally
$\Xi_{ra}^2\Xi_{rb}^2=0$ on every row for $a\neq b$, since at most one
coordinate is nonzero, proving \eqref{eq:gsv-cubature-cross}.
\end{proof}

\begin{remark}[Consequence at $d=18$ and empirical anchors]
\label{rem:gsv-defect-consequence}
The Austria SIR route ($d=18$, $N=1008$) must use $\Xi_c$ because a positive
replicated GenUT design is infeasible there
(Proposition~\ref{prop:gsv-genut-positivity}).  By
Proposition~\ref{prop:gsv-cubature-design-moments}, every reset injects
residual rows with per-axis kurtosis $18$ (Gaussian: $3$) and pairwise
co-kurtosis $0$ (Gaussian: $1$) into
$\widetilde y=y^{+}+\Xi_c B^{\top}$, whenever the injected scale
$\lVert B\rVert$ is non-negligible.  The higher-moment correction then
removes this artifact each step using noisy empirical targets and tuned
strengths, i.e.\ the recursion of
Proposition~\ref{prop:gsv-tangent-recursion} repeatedly processes a
designed-in shape error.  Empirically (descriptive, 16 fixed seeds, artifact
\texttt{austria\_sir\_pairwise\_moment\_genut\_score\_20260730}): the
diagonal-only score standard deviations were
$3435.6/1272.4/302.0$ for the three parameters, reduced by factors
$94.1/71.6/14.6$ by the pairwise co-moment correction --- consistent with the
mechanism, but the causal attribution to the design defect is a hypothesis
that Part~II's design swap is constructed to test.  No claim in this
subsection is a variance theorem.
\end{remark}

\section{Part II, Repair I: a fixed whitened Gaussian residual design}
\label{sec:gsv-design}

\begin{proposition}[Exact first- and second-moment identities]
\label{prop:gsv-whitened-design-identities}
Let $Z\in\mathbb{R}^{N\times d}$ with $N>d$, let
$\bar z=N^{-1}Z^{\top}\mathbf 1\in\mathbb{R}^{d}$,
$Z_c=Z-\mathbf 1\bar z^{\top}$, and suppose
$S=N^{-1}Z_c^{\top}Z_c$ is nonsingular with Cholesky factorization
$S=\widehat C\widehat C^{\top}$.  Define
$\Xi^{*}=Z_c\widehat C^{-\top}$.  Then, exactly,
\begin{equation}
\label{eq:gsv-whitened-identities}
 \frac1N(\Xi^{*})^{\top}\mathbf 1=0,
 \qquad
 \frac1N(\Xi^{*})^{\top}\Xi^{*}=I_d .
\end{equation}
If $Z$ is drawn once with a fixed recorded seed and never redrawn, then
$\Xi^{*}$ is a $\theta$-independent constant of the finite program.
\end{proposition}

\begin{proof}
$(\Xi^{*})^{\top}\mathbf 1
=\widehat C^{-1}Z_c^{\top}\mathbf 1
=\widehat C^{-1}\big(Z^{\top}\mathbf 1-N\bar z\big)=0$ by the definition of
$\bar z$.  Next,
$N^{-1}(\Xi^{*})^{\top}\Xi^{*}
=\widehat C^{-1}\big(N^{-1}Z_c^{\top}Z_c\big)\widehat C^{-\top}
=\widehat C^{-1}\widehat C\widehat C^{\top}\widehat C^{-\top}=I_d$.
The final statement holds because the construction involves no quantity that
depends on $\theta$.
\end{proof}

\begin{corollary}[The same-scalar score property is preserved]
\label{cor:gsv-design-tangent}
Replacing $\Xi_c$ by $\Xi^{*}$ leaves the chapter score proposition's
hypotheses intact: the residual design remains a fixed constant, so its
tangent is exactly zero and the recursive forward sensitivity remains the
total derivative of the executed scalar on the fixed branch.  Moreover the
divisibility constraint $N\equiv0\ (\mathrm{mod}\ 2d)$ required by the
replicated cubature construction is no longer needed.
\end{corollary}

\begin{proof}
The chapter score proposition fixes the residual design as part of the finite
route and differentiates all remaining dependencies; the only design property
it uses elsewhere is \eqref{eq:gsv-whitened-identities} (zero mean, identity
second moment), which Proposition~\ref{prop:gsv-whitened-design-identities}
supplies exactly.  $\theta$-independence gives $d\Xi^{*}/d\theta=0$, so the
implemented zero design-tangent is exact.  The replication argument that
forced $2d\mid N$ is not used by the construction of $\Xi^{*}$.
\end{proof}

\begin{proposition}[Raw moments of i.i.d.\ Gaussian columns]
\label{prop:gsv-iid-gaussian-moments}
Let $z_1,\dots,z_N$ and $w_1,\dots,w_N$ be independent standard normal
samples (two distinct coordinates of $Z$), and define
\begin{equation}
 m_2=\frac1N\sum_{r}z_r^{2},\quad
 m_4=\frac1N\sum_{r}z_r^{4},\quad
 m_2'=\frac1N\sum_{r}w_r^{2},\quad
 m_{22}=\frac1N\sum_{r}z_r^{2}w_r^{2}.
\end{equation}
Then
\begin{equation}
\label{eq:gsv-iid-moments}
\begin{aligned}
 &\mathbb{E}m_2=1,\quad \operatorname{Var}m_2=\tfrac{2}{N},\qquad
 \mathbb{E}m_4=3,\quad \operatorname{Var}m_4=\tfrac{96}{N},\qquad
 \operatorname{Cov}(m_2,m_4)=\tfrac{12}{N},\\
 &\mathbb{E}m_{22}=1,\quad \operatorname{Var}m_{22}=\tfrac{8}{N},\qquad
 \operatorname{Cov}(m_{22},m_2)=\tfrac{2}{N},\qquad
 \operatorname{Cov}(m_2,m_2')=0 .
\end{aligned}
\end{equation}
\end{proposition}

\begin{proof}
Standard normal moments: $\mathbb{E}z^{2}=1$, $\mathbb{E}z^{4}=3$,
$\mathbb{E}z^{6}=15$, $\mathbb{E}z^{8}=105$.  For i.i.d.\ averages,
$\operatorname{Var}(N^{-1}\sum_rY_r)=\operatorname{Var}(Y)/N$ and
$\operatorname{Cov}$ likewise, because distinct $r$ are independent.  Thus
$\operatorname{Var}m_2=(\mathbb{E}z^{4}-1)/N=2/N$;
$\operatorname{Var}m_4=(\mathbb{E}z^{8}-9)/N=96/N$;
$\operatorname{Cov}(m_2,m_4)=(\mathbb{E}z^{6}-\mathbb{E}z^{2}\mathbb{E}z^{4})/N
=(15-3)/N=12/N$.
With $z\perp w$:
$\mathbb{E}m_{22}=\mathbb{E}z^{2}\mathbb{E}w^{2}=1$;
$\operatorname{Var}m_{22}
=(\mathbb{E}z^{4}\mathbb{E}w^{4}-1)/N=(9-1)/N=8/N$;
$\operatorname{Cov}(m_{22},m_2)
=(\mathbb{E}z^{4}\mathbb{E}w^{2}-1)/N=(3-1)/N=2/N$;
and $\operatorname{Cov}(m_2,m_2')=0$ by independence of the coordinates.
\end{proof}

\begin{proposition}[Delta-method asymptotics for the studentized shape statistics]
\label{prop:gsv-delta-method}
In the setting of Proposition~\ref{prop:gsv-iid-gaussian-moments} define the
studentized kurtosis and pair co-kurtosis
\begin{equation}
 \kappa_N=\frac{m_4}{m_2^{2}},
 \qquad
 c_N=\frac{m_{22}}{m_2\,m_2'} .
\end{equation}
Then, as $N\to\infty$,
\begin{equation}
\label{eq:gsv-delta}
 \sqrt N\,(\kappa_N-3)\;\Rightarrow\;\mathcal N(0,\,24),
 \qquad
 \sqrt N\,(c_N-1)\;\Rightarrow\;\mathcal N(0,\,4).
\end{equation}
At $N=1008$ the corresponding studentized standard deviations are approximately
$\sqrt{24/N}\approx0.154$ and $\sqrt{4/N}\approx0.063$.  The raw
$m_{22}$ statistic has $\sqrt{8/N}$ standard deviation; that is a different
quantity and should not be used as the whitened co-kurtosis fluctuation.
\end{proposition}

\begin{proof}
By the multivariate CLT and \eqref{eq:gsv-iid-moments},
$\sqrt N\big((m_4,m_2)-(3,1)\big)\Rightarrow\mathcal N(0,\Sigma_1)$ with
$\Sigma_1=\begin{pmatrix}96&12\\12&2\end{pmatrix}$.
For $g(m_4,m_2)=m_4/m_2^{2}$ the gradient at $(3,1)$ is
$\nabla g=(1,\,-6)$ since $\partial g/\partial m_4=m_2^{-2}=1$ and
$\partial g/\partial m_2=-2m_4m_2^{-3}=-6$.  The delta method gives
asymptotic variance
\begin{equation}
\label{eq:gsv-delta-kurt-arithmetic}
 \nabla g^{\top}\Sigma_1\nabla g
 =96\cdot1-12\cdot6-6\cdot12+36\cdot2
 =96-72-72+72=24 .
\end{equation}
Similarly
$\sqrt N\big((m_{22},m_2,m_2')-(1,1,1)\big)\Rightarrow\mathcal N(0,\Sigma_2)$
with
$\Sigma_2=\begin{pmatrix}8&2&2\\2&2&0\\2&0&2\end{pmatrix}$,
and for $h(m_{22},m_2,m_2')=m_{22}/(m_2m_2')$ the gradient at $(1,1,1)$ is
$(1,\,-1,\,-1)$, giving
\begin{equation}
\label{eq:gsv-delta-cokurt-arithmetic}
 \nabla h^{\top}\Sigma_2\nabla h
 =8+2+2-2\cdot2-2\cdot2+2\cdot0=4 .
\end{equation}
\end{proof}

\begin{remark}[Whitening versus studentization: a stated gap]
\label{rem:gsv-whitening-gap}
Proposition~\ref{prop:gsv-delta-method} treats per-coordinate studentization.
The design $\Xi^{*}$ of
Proposition~\ref{prop:gsv-whitened-design-identities} applies full
$d$-dimensional mean-and-Cholesky whitening, which fixes \emph{all} first and
second moments exactly, including cross terms.  The whitened fourth-order
statistics agree with the studentized ones to first order (both are smooth
functions of the same sample moments with the same linearization at the
Gaussian point), so \eqref{eq:gsv-delta} is the first-order description used
here; the $O_p(N^{-1})$ correction induced by the off-diagonal
whitening constraints is \emph{not} derived in this note and is recorded as
a gap, not a claim.  For the intended use --- replacing exact defects of size
$d-3=15$ and $-1$ by fluctuations of size $\approx0.15$ and $\approx0.06$ ---
first-order accuracy is ample.
\end{remark}

\begin{proposition}[Why a positive GenUT design cannot do this for $d\ge4$]
\label{prop:gsv-genut-positivity}
The Gaussian-moment GenUT axis construction with standardized skewness
$s_a=0$ and kurtosis $k_a=3$ yields, per axis,
$u_a=v_a=\sqrt3$ and noncentral weights $b_a=c_a=\tfrac16$, hence central
weight
\begin{equation}
\label{eq:gsv-genut-central-weight}
 w_0=1-\sum_{a=1}^{d}(b_a+c_a)=1-\frac{d}{3},
\end{equation}
which is negative for every $d\ge4$.  Consequently no positive equal-weight
replication of the $2d{+}1$-point Gaussian GenUT rule exists for $d\ge4$,
while $\Xi^{*}$ is positive (equal-weight by construction) at every $d$ and
$N>d$.
\end{proposition}

\begin{proof}
The axis equations of the GenUT construction with $s_a=0$, $k_a=3$ give
$u_a=\tfrac{-s_a+\sqrt{4k_a-3s_a^{2}}}{2}=\tfrac{\sqrt{12}}{2}=\sqrt3$,
$v_a=u_a+s_a=\sqrt3$, and
$c_a=\frac{1}{v_a(u_a+v_a)}=\frac{1}{\sqrt3\cdot2\sqrt3}=\frac16$,
$b_a=\frac{1}{u_a(u_a+v_a)}=\frac16$.  Summing over the $d$ axes gives
\eqref{eq:gsv-genut-central-weight}; $1-d/3<0$ iff $d>3$.  A rule with a
negative weight cannot be represented by nonnegative equal-weight
replication, which requires every atom mass to be a nonnegative multiple of
$1/N$.
\end{proof}

\section{Part II, Repair IV: frozen low-rank projected cumulant matching}
\label{sec:gsv-projected-cumulants}

\subsection{Source boundary and target quantity}

This section defines a new opt-in correction.  It is motivated by the general
low-rank principle used by tensor-train methods, but it is not the moment
algorithm of Zhao and Cui.  Zhao and Cui approximate a sequential density by
a TT or squared-TT representation, marginalize its cores, and construct
conditional Knothe--Rosenblatt maps (their equations (12)--(16), Algorithms
1--2, and Proposition~2).  Their author implementation performs the adjacent
fit in
\texttt{third\_party/audit/zhao\_cui\_tensor\_ssm\_p10/source/models/full\_sol.m},
lines 21--136, and squared-TT marginal contraction in
\texttt{third\_party/audit/zhao\_cui\_tensor\_ssm\_p10/source/deep-tensor.dev/src/@TTSIRT/marginalise.m},
lines 25--85.  Neither inspected source defines the projected cumulant update
below.  The classification is therefore
\texttt{extension\_or\_invention}, not \texttt{source\_faithful}.

Let $z_i\in\mathbb R^d$ denote the weighted source cloud after exact weighted
whitening and let $q_i\in\mathbb R^d$ denote the current equal-weight reset
cloud after exact uniform whitening and any declared diagonal correction.
Thus
\begin{equation}
 \sum_i\alpha_i z_i=0,\quad
 \sum_i\alpha_i z_i z_i^\top=I_d,\qquad
 \frac1N\sum_iq_i=0,\quad
 \frac1N\sum_iq_iq_i^\top=I_d.
 \label{eq:gsv-lr-whitened-clouds}
\end{equation}
Fix an orthonormal matrix $U\in\mathbb R^{d\times r}$, $U^\top U=I_r$,
selected offline and frozen for the complete execution scope.  Define
$y_i=U^\top z_i$ and $s_i=U^\top q_i$.  The projected third central moments
and fourth cumulants are
\begin{align}
 T_3^z&=\sum_i\alpha_i y_i^{\otimes3}, &
 T_3^q&=\frac1N\sum_i s_i^{\otimes3},
 \label{eq:gsv-lr-third}\\
 K_4^z&=\sum_i\alpha_i y_i^{\otimes4}-\mathcal W_r, &
 K_4^q&=\frac1N\sum_i s_i^{\otimes4}-\mathcal W_r.
 \label{eq:gsv-lr-fourth}
\end{align}
where
$(\mathcal W_r)_{abcd}=\delta_{ab}\delta_{cd}+
\delta_{ac}\delta_{bd}+\delta_{ad}\delta_{bc}$ is the Gaussian Wick tensor.
Because both projected covariances equal $I_r$, the Wick terms cancel in the
residual.  Set
\begin{equation}
 R_3=T_3^z-T_3^q,\qquad R_4=K_4^z-K_4^q
 =\sum_i\alpha_i y_i^{\otimes4}-\frac1N\sum_i s_i^{\otimes4}.
 \label{eq:gsv-lr-residuals}
\end{equation}
Unlike the pairwise statistics $E[z_a^2z_b]$ and $E[z_a^2z_b^2]$, these are
the complete order-three and order-four tensors inside the selected
$r$-dimensional subspace.  They include interactions among three or four
distinct projected coordinates.

\subsection{One bounded correction step}

For each current projected row $s_i$, define the residual-gradient direction
\begin{equation}
 (g_i)_a=
 3\sum_{b,c=1}^r(R_3)_{abc}(s_i)_b(s_i)_c
 +4\sum_{b,c,e=1}^r(R_4)_{abce}(s_i)_b(s_i)_c(s_i)_e.
 \label{eq:gsv-lr-projected-direction}
\end{equation}
The lifted direction is $D_i=Ug_i$ in column-vector notation.  To preserve
the affine constraints to first order, remove its uniform mean and its
symmetric covariance tangent:
\begin{align}
 \bar D&=N^{-1}\sum_iD_i,\qquad C_i=D_i-\bar D,\
 A&=\operatorname{sym}\!\left(N^{-1}\sum_i q_iC_i^\top\right),
 &P_i&=C_i-Aq_i .
 \label{eq:gsv-lr-affine-projection}
\end{align}
With a declared floor $\delta>0$, bounded strength $\rho\ge0$, and
\begin{equation}
 a=\left(N^{-1}\sum_i\lVert P_i\rVert_2^2+\delta\right)^{1/2},
 \qquad \widetilde q_i=q_i+\rho P_i/a,
 \label{eq:gsv-lr-bounded-update}
\end{equation}
apply exact uniform mean/covariance restandardization to
$\{\widetilde q_i\}$.  Repeat a fixed number of times.  Finally restore the
weighted source mean and Cholesky covariance exactly, as in the existing
higher-moment Contract~E composition.  Rank, basis, step count, strength, and
floor are fixed route controls; no runtime rank selection is permitted.

\begin{proposition}[Direction is the projected moment-residual gradient]
\label{prop:gsv-lr-gradient}
Let
\begin{equation}
 \mathcal J(q)=\tfrac12\lVert T_3^z-T_3^q\rVert_F^2
 +\tfrac12\lVert K_4^z-K_4^q\rVert_F^2,
 \label{eq:gsv-lr-objective}
\end{equation}
with the target tensors held fixed.  Before the affine projection and common
positive scaling, the negative gradient of $\mathcal J$ with respect to the
projected row $s_i$ is $N^{-1}g_i$, with $g_i$ given by
\eqref{eq:gsv-lr-projected-direction}.
\end{proposition}

\begin{proof}
For a perturbation $h_i\in\mathbb R^r$,
$d(s_i^{\otimes k})[h_i]$ is the sum of the $k$ tensors obtained by replacing
one factor $s_i$ by $h_i$.  Both $R_3$ and $R_4$ are symmetric, so contracting
these sums with $R_k$ gives respectively
$3\langle R_3,h_i\otimes s_i\otimes s_i\rangle/N$ and
$4\langle R_4,h_i\otimes s_i^{\otimes3}\rangle/N$.
Since $dR_k=-dT_k^q$, differentiating
$\tfrac12\lVert R_k\rVert_F^2$ supplies the minus sign.  Reading off the
coefficient of $h_i$ gives $-g_i/N$.
\end{proof}

\begin{proposition}[Exact affine restoration]
\label{prop:gsv-lr-affine-restoration}
Suppose the covariance of $\{\widetilde q_i\}$ in
\eqref{eq:gsv-lr-bounded-update} is positive definite.  Uniform
mean/Cholesky restandardization gives an output with exactly zero empirical
mean and identity empirical covariance in exact arithmetic.  Subsequent
restoration by the source mean and Cholesky factor gives exactly the weighted
source mean and covariance.
\end{proposition}

\begin{proof}
This is the same algebra as
Proposition~\ref{prop:gsv-whitened-design-identities}: subtract the uniform
mean and right-solve by the Cholesky factor of the uniform covariance.  The
affine restoration then maps zero and $I_d$ to the source mean and source
covariance.  The projection in \eqref{eq:gsv-lr-affine-projection} reduces
the size of the restandardization but is not used to claim exactness; the
final Cholesky operation supplies the exact identities.
\end{proof}

\subsection{Total tangent of the same finite program}

The basis $U$ has zero tangent because it is learned on disjoint data and
frozen before the claim run.  For any parameter direction, dots denote total
derivatives.  Projected rows satisfy
$\dot y_i=U^\top\dot z_i$ and $\dot s_i=U^\top\dot q_i$.  For $k=3,4$,
\begin{align}
 \dot T_k^z
 &=\sum_i\dot\alpha_i y_i^{\otimes k}
 +\sum_i\alpha_i\sum_{j=1}^k
 y_i^{\otimes(j-1)}\otimes\dot y_i\otimes
 y_i^{\otimes(k-j)},
 \label{eq:gsv-lr-target-jvp}\\
 \dot T_k^q
 &=\frac1N\sum_i\sum_{j=1}^k
 s_i^{\otimes(j-1)}\otimes\dot s_i\otimes
 s_i^{\otimes(k-j)},
 \\
 \dot R_k=\dot T_k^z-\dot T_k^q .
 \label{eq:gsv-lr-current-jvp}
\end{align}
Differentiating \eqref{eq:gsv-lr-projected-direction} gives, in index form,
\begin{align}
 (\dot g_i)_a={}&3\sum_{bc}
 \left[(\dot R_3)_{abc}s_{ib}s_{ic}
 +(R_3)_{abc}(\dot s_{ib}s_{ic}+s_{ib}\dot s_{ic})\right] \\
 &+4\sum_{bce}\left[(\dot R_4)_{abce}s_{ib}s_{ic}s_{ie}
 +(R_4)_{abce}(\dot s_{ib}s_{ic}s_{ie}
 +s_{ib}\dot s_{ic}s_{ie}+s_{ib}s_{ic}\dot s_{ie})\right].
 \label{eq:gsv-lr-direction-jvp}
\end{align}
The JVP of \eqref{eq:gsv-lr-affine-projection} uses the product rule through
$\bar D$, $A$, and $P_i$.  The scale derivative is
\begin{equation}
 \dot a=\frac{N^{-1}\sum_iP_i^\top\dot P_i}{a},qquad
 \frac{d}{d\theta}(P_i/a)=\dot P_i/a-P_i\dot a/a^2.
 \label{eq:gsv-lr-scale-jvp}
\end{equation}
The existing Cholesky JVP then differentiates both restandardization and
source affine restoration.  Equations
\eqref{eq:gsv-lr-target-jvp}--\eqref{eq:gsv-lr-scale-jvp} contain target
weight, target-cloud, current-cloud, projection, normalization, and
restoration dependence.  Omitting any of these terms is a partial derivative
and is wrong relative to the same-scalar score claim.

\begin{corollary}[Same-scalar derivative on the fixed branch]
\label{cor:gsv-lr-same-scalar}
If $U$, $r$, and all correction controls are frozen independently of the
claim parameter and observations, and every Cholesky/validity branch remains
fixed, composing the JVP above with the existing total OT and Contract-E JVP
gives the total derivative of the executed finite scalar.
\end{corollary}

\begin{proof}
The projected correction is a finite composition of polynomial contractions,
linear projections, positive square roots, and Cholesky solves.  Equations
\eqref{eq:gsv-lr-target-jvp}--\eqref{eq:gsv-lr-scale-jvp} are its chain rule;
the frozen basis contributes no omitted tangent.  Composition with the
already total upstream JVP proves the statement.
\end{proof}

\subsection{Offline subspace construction and stability diagnostic}

Whitening makes covariance PCA uninformative: every direction has unit
second moment.  The basis must therefore be learned from higher-moment
\emph{residuals}.  Use fixed recorded sketch directions
$a_1,\ldots,a_S\in\mathbb R^d$ and, for each disjoint calibration cloud pair,
form mode-one residual sketches
\begin{align}
 G_{3,s}&=\sum_i\alpha_i z_i(a_s^\top z_i)^2
 -\frac1N\sum_iq_i(a_s^\top q_i)^2,\\
 G_{4,s}&=\sum_i\alpha_i z_i(a_s^\top z_i)^3
 -\frac1N\sum_iq_i(a_s^\top q_i)^3.
 \label{eq:gsv-lr-sketches}
\end{align}
Accumulate the positive semidefinite mode score
\begin{equation}
 H=\sum_{\text{calibration clouds}}\sum_{s=1}^S
 \left(G_{3,s}G_{3,s}^\top+G_{4,s}G_{4,s}^\top\right).
 \label{eq:gsv-lr-mode-score}
\end{equation}
The nested rank-eight basis consists of the leading eight eigenvectors of
$H$; ranks four and six take its first four and six columns.  The sketch seed,
calibration observations, particle seeds, scaling between orders, and
eigenvalue tie/sign convention are part of the issued tuning artifact.
Fit another basis on untouched validation clouds and report principal angles,
and report the validation explained score
$\operatorname{tr}(U_r^\top H_{\rm val}U_r)/
\operatorname{tr}(H_{\rm val})$.  These diagnose whether low-rank structure
is reproducible rather than particle noise.  The validation basis is never
used in the claim computation.

\subsection{Cost, implementation contract, and non-claims}

For one correction step, projection costs $O(Ndr)$, forming the complete
projected moments costs $O(Nr^4)$, their contractions cost $O(Nr^4)$, and
lifting/projecting the update costs $O(Ndr+Nd^2)$.  Storage is
$O(N(d+r)+r^4)$ beyond the existing tangent arrays.  At $d=18$, ranks
$r=4,6,8$ require respectively $4^4=256$, $6^4=1296$, and $8^4=4096$
fourth-order entries, versus $18^4=104{,}976$ for a dense ambient tensor.
The runtime path must use TensorFlow/XLA contractions and must never
materialize ambient $d^3$ or $d^4$ tensors.

This construction changes the finite reset map and therefore generally
changes both value and same-scalar score.  Low residual energy does not prove
lower likelihood or score bias.  A stable low-rank basis does not prove that
ranks remain bounded with dimension or horizon.  Empirical value stability,
score dispersion, tangent growth, derivative parity, and downstream HMC
remain separate gates.  In particular, the absence of an exact nonlinear
Austria score oracle prevents an absolute score-accuracy claim.

\section{Part II, Repair II: surrogate-force HMC with the exact executed energy}
\label{sec:gsv-hmc}

Throughout this section $U:\mathbb{R}^{P}\to\mathbb{R}$ is the executed
finite potential $U(\theta)=-\widehat L^{N}(\theta)$ with all seeds frozen
(a deterministic function), $M\succ0$ a fixed mass matrix,
$H(\theta,p)=U(\theta)+\tfrac12p^{\top}M^{-1}p$, and
$F:\mathbb{R}^{P}\to\mathbb{R}^{P}$ an \emph{arbitrary} measurable
deterministic map (``force''), not assumed to be a gradient.

\begin{definition}[Force-field leapfrog]
\label{def:gsv-leapfrog}
For step size $\epsilon>0$, one leapfrog step
$\Psi:(\theta_0,p_0)\mapsto(\theta_1,p_1)$ is
\begin{equation}
\label{eq:gsv-leapfrog}
 p_{1/2}=p_0+\tfrac{\epsilon}{2}F(\theta_0),
 \qquad
 \theta_1=\theta_0+\epsilon M^{-1}p_{1/2},
 \qquad
 p_1=p_{1/2}+\tfrac{\epsilon}{2}F(\theta_1).
\end{equation}
\end{definition}

\begin{proposition}[Measure preservation for any force]
\label{prop:gsv-leapfrog-volume}
Each substep of \eqref{eq:gsv-leapfrog}, and hence $\Psi$ and any iterate
$\Psi^{L}$, preserves Lebesgue measure on
$\mathbb{R}^{P}\times\mathbb{R}^{P}$, for every measurable $F$.
\end{proposition}

\begin{proof}
The first and third substeps have the form
$(\theta,p)\mapsto(\theta,\,p+a(\theta))$; for each fixed $\theta$ this is a
translation of $p$, so by Fubini's theorem it preserves the product measure.
The middle substep $(\theta,p)\mapsto(\theta+\epsilon M^{-1}p,\,p)$ is, for
each fixed $p$, a translation of $\theta$, and preserves measure by the same
argument.  A composition of measure-preserving bijections is
measure-preserving.
\end{proof}

\begin{proposition}[Reversibility for any force]
\label{prop:gsv-leapfrog-reversible}
Let $\mathcal F(\theta,p)=(\theta,-p)$ be the momentum flip.  Then for every
$F$ and every $L\ge1$,
\begin{equation}
\label{eq:gsv-involution}
 \big(\mathcal F\circ\Psi^{L}\big)\circ\big(\mathcal F\circ\Psi^{L}\big)
 =\mathrm{id},
\end{equation}
i.e.\ $S=\mathcal F\circ\Psi^{L}$ is an involution.
\end{proposition}

\begin{proof}
First take $L=1$.  Apply $\Psi$ to $\mathcal F(\theta_1,p_1)=(\theta_1,-p_1)$
using \eqref{eq:gsv-leapfrog}: the half-kick gives
$-p_1+\tfrac{\epsilon}{2}F(\theta_1)
=-\big(p_{1/2}+\tfrac{\epsilon}{2}F(\theta_1)\big)
 +\tfrac{\epsilon}{2}F(\theta_1)=-p_{1/2}$;
the drift gives
$\theta_1+\epsilon M^{-1}(-p_{1/2})=\theta_0$;
the final half-kick gives
$-p_{1/2}+\tfrac{\epsilon}{2}F(\theta_0)
=-\big(p_0+\tfrac{\epsilon}{2}F(\theta_0)\big)
 +\tfrac{\epsilon}{2}F(\theta_0)=-p_0$.
Hence $\Psi\circ\mathcal F\circ\Psi=\mathcal F$, equivalently
$\mathcal F\circ\Psi\circ\mathcal F\circ\Psi=\mathrm{id}$ (using
$\mathcal F\circ\mathcal F=\mathrm{id}$).  For $L>1$, iterate:
$\Psi^{L}\circ\mathcal F\circ\Psi^{L}
=\Psi^{L-1}\circ(\Psi\circ\mathcal F\circ\Psi)\circ\Psi^{L-1}
=\Psi^{L-1}\circ\mathcal F\circ\Psi^{L-1}=\cdots=\mathcal F$.
Note that no palindromic step arrangement is needed because $F$ depends only
on $\theta$, so the same substep sequence serves both directions.
\end{proof}

\begin{proposition}[Invariance of Metropolis--Hastings with an involutive
measure-preserving proposal]
\label{prop:gsv-mh-invariance}
Let $\pi$ be a probability density, positive on
$\mathbb{R}^{P}\times\mathbb{R}^{P}$, and let $S$ be a measurable,
Lebesgue-measure-preserving involution.  Define the kernel: at state $z$,
propose $z'=S(z)$ and accept with probability
$a(z)=\min\{1,\pi(S(z))/\pi(z)\}$, else stay.  Then the kernel is
$\pi$-reversible, hence leaves $\pi$ invariant.
\end{proposition}

\begin{proof}
Let $m(z)=\pi(z)\,a(z)=\min\{\pi(z),\pi(S(z))\}$.  Since $S$ is an
involution, $m(S(z))=\min\{\pi(S(z)),\pi(z)\}=m(z)$.  For bounded measurable
$f,g$, the accepted flow satisfies, by the change of variables $z=S(w)$
(measure-preserving bijection),
\begin{equation}
\label{eq:gsv-detailed-balance}
 \int f(z)\,g(S(z))\,m(z)\,dz
 =\int f(S(w))\,g(w)\,m(S(w))\,dw
 =\int g(w)\,f(S(w))\,m(w)\,dw ,
\end{equation}
which is the same expression with $f$ and $g$ interchanged.  The rejected
flow $\int f(z)g(z)\,(\pi(z)-m(z))\,dz$ is symmetric in $(f,g)$ trivially.
Summing gives
$\mathbb{E}_{\pi}[f(z)g(z^{+})]=\mathbb{E}_{\pi}[g(z)f(z^{+})]$ for the
chain step $z\to z^{+}$, i.e.\ $\pi$-reversibility.
\end{proof}

\begin{corollary}[Surrogate-force HMC targets the exact executed scalar]
\label{cor:gsv-surrogate-force}
Consider the algorithm: (i) refresh $p\sim\mathcal N(0,M)$; (ii) propose
$S(\theta,p)$ with $S=\mathcal F\circ\Psi^{L}$ built from \emph{any}
deterministic momentum-independent force $F$; (iii) accept with probability
$\min\{1,e^{H(\theta,p)-H(S(\theta,p))}\}$ where $H$ uses the exact executed
potential $U=-\widehat L^{N}$.  Every step leaves
$\pi(\theta,p)\propto e^{-H(\theta,p)}$ invariant, hence the $\theta$-marginal
$\propto e^{-U(\theta)}$ is preserved, for any quality of $F$.  The choice of
$F$ affects only the acceptance probability and mixing.
\end{corollary}

\begin{proof}
Step (i) is a Gibbs update of the $p$-coordinate from its exact conditional
under $\pi$ (which is $\mathcal N(0,M)$ independent of $\theta$), hence
$\pi$-invariant.  Steps (ii)--(iii) form the kernel of
Proposition~\ref{prop:gsv-mh-invariance} with the involutive
measure-preserving $S$ supplied by
Propositions~\ref{prop:gsv-leapfrog-volume} and
\ref{prop:gsv-leapfrog-reversible}, hence are $\pi$-reversible.  A
composition of $\pi$-invariant kernels is $\pi$-invariant.  Since $F$ enters
only through the proposal path, not through $\pi$ or $H$, its quality cannot
affect invariance.
\end{proof}

\begin{remark}[Requirements, and what this buys]
\label{rem:gsv-force-requirements}
$F$ must be a deterministic function of $\theta$ alone: all Monte Carlo seeds
inside $F$ must be frozen, and $F$ may not depend on momentum or on the
trajectory's history.  Under those conditions
Corollary~\ref{cor:gsv-surrogate-force} licenses using, as leapfrog force,
(i) a damped or clipped variant of the recursive score (e.g.\ computed with
larger $\lambda,\delta$ than the value program), (ii) the Fisher-identity
estimator of Section~\ref{sec:gsv-fisher} with its own frozen seeds, or
(iii) a deterministic teacher or sigma-point score --- while the acceptance
step keeps the \emph{exact} same-scalar energy, so correctness never depends
on score variance again.  The current repository HMC wrapper passes a single
custom-gradient target (one adapter call returns value and score), so this
scheme is an implementable proposal, not current behavior.  No claim is made
about acceptance rates or mixing; those are empirical questions for a
planned experiment.
\end{remark}

\section{Part II, Repair III: Fisher-identity estimators without transport
derivatives}
\label{sec:gsv-fisher}

This section concerns Object~B of Definition~\ref{def:gsv-objects}.  Consider
the state-space model with initial density $\mu_\theta(x_0)$, transition
densities $f_\theta(x_t\mid x_{t-1})$, and observation densities
$g_\theta(y_t\mid x_t)$, all positive and continuously differentiable in
$\theta$.

\begin{assumption}[Interchange]
\label{ass:gsv-interchange}
For $\theta$ in an open set, the joint density
$p_\theta(x_{0:T},y_{1:T})
=\mu_\theta(x_0)\prod_{t=1}^{T}f_\theta(x_t\mid x_{t-1})
 \,g_\theta(y_t\mid x_t)$
is differentiable in $\theta$ with
$\lvert\partial_{\theta_p}p_\theta(x,y)\rvert\le G(x)$ for an integrable
envelope $G$, locally uniformly in $\theta$, so that
$\partial_{\theta}\!\int p_\theta\,dx=\int\partial_\theta p_\theta\,dx$.
\end{assumption}

\begin{proposition}[Fisher identity and additive decomposition]
\label{prop:gsv-fisher-identity}
Under Assumption~\ref{ass:gsv-interchange}, with
$p_\theta(y_{1:T})=\int p_\theta(x_{0:T},y_{1:T})\,dx_{0:T}>0$,
\begin{equation}
\label{eq:gsv-fisher}
 \nabla_\theta\log p_\theta(y_{1:T})
 =\mathbb{E}\!\left[\nabla_\theta\log p_\theta(x_{0:T},y_{1:T})
 \,\middle|\,y_{1:T}\right]
 =\mathbb{E}\!\left[\sum_{t=0}^{T}s_t\,\middle|\,y_{1:T}\right],
\end{equation}
where $s_0=\nabla_\theta\log\mu_\theta(x_0)$ and, for $t\ge1$,
\begin{equation}
\label{eq:gsv-additive}
 s_t=s_t(x_{t-1},x_t)
 =\nabla_\theta\log f_\theta(x_t\mid x_{t-1})
 +\nabla_\theta\log g_\theta(y_t\mid x_t).
\end{equation}
No derivative of any resampling, transport, or reset operation appears in
\eqref{eq:gsv-fisher}--\eqref{eq:gsv-additive}.
\end{proposition}

\begin{proof}
By Assumption~\ref{ass:gsv-interchange},
$\nabla\log p(y)
=\frac{\int\nabla p_\theta(x,y)\,dx}{p_\theta(y)}
=\int\nabla\log p_\theta(x,y)\,
 \frac{p_\theta(x,y)}{p_\theta(y)}\,dx
=\mathbb{E}[\nabla\log p_\theta(x,y)\mid y]$,
using $\nabla p=p\,\nabla\log p$ and the definition of the smoothing density
$p_\theta(x\mid y)=p_\theta(x,y)/p_\theta(y)$.  Taking the logarithm of the
displayed factorization of $p_\theta(x,y)$ makes
$\nabla\log p_\theta(x,y)=\sum_{t}s_t$, which is
\eqref{eq:gsv-additive}.  The final sentence holds because the identity is a
property of the model densities, independent of any algorithm used to
approximate the conditional expectation.
\end{proof}

\begin{lemma}[Backward factorization]
\label{lem:gsv-backward-kernel}
For $t\ge1$, with all densities positive:
\begin{align}
 p(x_{t-1}\mid x_t,y_{1:t})
 &=p(x_{t-1}\mid x_t,y_{1:t-1})
 =\frac{p(x_{t-1}\mid y_{1:t-1})\,f_\theta(x_t\mid x_{t-1})}
        {\int p(x'\mid y_{1:t-1})\,f_\theta(x_t\mid x')\,dx'} ,
 \label{eq:gsv-backward-kernel}\\
 p(x_{0:t-2}\mid x_{t-1},x_t,y_{1:t})
 &=p(x_{0:t-2}\mid x_{t-1},y_{1:t-1}).
 \label{eq:gsv-past-independence}
\end{align}
\end{lemma}

\begin{proof}
For \eqref{eq:gsv-backward-kernel}: $y_t$ is conditionally independent of
$x_{t-1}$ given $x_t$ (the observation density depends on $x_t$ only), which
removes $y_t$ from the conditioning; then Bayes' rule with the Markov
transition gives
$p(x_{t-1}\mid x_t,y_{1:t-1})
=\frac{p(x_{t-1},x_t\mid y_{1:t-1})}{p(x_t\mid y_{1:t-1})}
=\frac{p(x_{t-1}\mid y_{1:t-1})f_\theta(x_t\mid x_{t-1})}
       {p(x_t\mid y_{1:t-1})}$,
whose denominator is the displayed normalizer.  For
\eqref{eq:gsv-past-independence}: given $(x_{t-1},y_{1:t-1})$, the pair
$(x_t,y_t)$ depends on the past only through $x_{t-1}$ by the Markov and
observation structure, so conditioning on it does not change the law of
$x_{0:t-2}$.
\end{proof}

\begin{proposition}[Backward recursion for the smoothed additive functional]
\label{prop:gsv-backward-recursion}
Define $T_t(x_t)=\mathbb{E}\big[\sum_{s=0}^{t}s_s\,\big|\,x_t,y_{1:t}\big]$
with $T_0(x_0)=s_0(x_0)$.  Then for $t\ge1$,
\begin{equation}
\label{eq:gsv-backward-recursion}
 T_t(x_t)
 =\int p(x_{t-1}\mid x_t,y_{1:t})
 \Big[T_{t-1}(x_{t-1})+s_t(x_{t-1},x_t)\Big]\,dx_{t-1},
\end{equation}
and the exact score is
$\nabla_\theta\log p_\theta(y_{1:T})
=\mathbb{E}\big[T_T(x_T)\,\big|\,y_{1:T}\big]$,
an expectation under the filter marginal at time $T$.
\end{proposition}

\begin{proof}
By the tower property, conditioning first on $(x_{t-1},x_t,y_{1:t})$:
\[
 T_t(x_t)
 =\mathbb{E}\Big[\;
 \mathbb{E}\big[\textstyle\sum_{s\le t-1}s_s\,\big|\,x_{t-1},x_t,y_{1:t}\big]
 +s_t(x_{t-1},x_t)
 \;\Big|\;x_t,y_{1:t}\Big].
\]
By \eqref{eq:gsv-past-independence} the inner conditional expectation equals
$\mathbb{E}[\sum_{s\le t-1}s_s\mid x_{t-1},y_{1:t-1}]=T_{t-1}(x_{t-1})$, and
the outer expectation integrates $x_{t-1}$ against
\eqref{eq:gsv-backward-kernel}, giving
\eqref{eq:gsv-backward-recursion}.  Finally, by the tower property again,
$\mathbb{E}[T_T(x_T)\mid y_{1:T}]
=\mathbb{E}[\sum_{t\le T}s_t\mid y_{1:T}]$, which is the score by
Proposition~\ref{prop:gsv-fisher-identity}.
\end{proof}

\begin{definition}[Particle forward-smoothing estimator]
\label{def:gsv-pds}
Given any particle approximation
$\sum_j\beta_{t-1,j}\,\delta_{z_{t-1,j}}$ of $p(x_{t-1}\mid y_{1:t-1})$ and
forward particles $x_{t,i}^{-}$, plug it into
\eqref{eq:gsv-backward-kernel}--\eqref{eq:gsv-backward-recursion}:
\begin{equation}
\label{eq:gsv-pds}
 W_t(i,j)
 =\frac{\beta_{t-1,j}\,f_\theta\!\big(x_{t,i}^{-}\mid z_{t-1,j}\big)}
        {\sum_k\beta_{t-1,k}\,f_\theta\!\big(x_{t,i}^{-}\mid z_{t-1,k}\big)},
 \qquad
 \widehat T_{t,i}
 =\sum_j W_t(i,j)\Big[\widehat T_{t-1,j}
 +s_t\big(z_{t-1,j},x_{t,i}^{-}\big)\Big],
\end{equation}
and report $\widehat S_T=\sum_i\alpha_{T,i}\widehat T_{T,i}$.
\end{definition}

\begin{remark}[Cost, variance status, and route-compatibility boundary]
\label{rem:gsv-pds-boundary}
(i) The recursion \eqref{eq:gsv-pds} costs $O(N^{2})$ transition-density
evaluations per step ($\approx10^{6}$ Gaussian evaluations at $N=1008$),
batched trivially on GPU; every current model adapter has an evaluable
additive-Gaussian transition density and closed-form
$\nabla_\theta\log f_\theta$ reusing the existing transition tangents.
(ii) The estimator targets Object~B; used as an HMC gradient it must go
through Corollary~\ref{cor:gsv-surrogate-force}, never through the accept
step.  (iii) The literature result that the variance of \eqref{eq:gsv-pds}
grows linearly in $T$, versus quadratically for the ancestral path
estimator, is \emph{cited, not re-proved here} (Poyiadjis, Doucet and Singh,
Biometrika 2011); the local-copy policy requires storing and reading that
paper before an implementation decision.  (iv) On the OT/moment-reset route
the forward particles are propagated from a reset cloud that is a changed
object relative to categorical resampling; \eqref{eq:gsv-pds} then inherits
that approximation through the atom locations.  As a bias \emph{reference
oracle} the estimator should therefore be run on a plain
multinomial-resampling particle filter arm, where its only error is
quantifiable Monte Carlo error.
\end{remark}

\section{A branch warning for ESS-triggered resets}
\label{sec:gsv-branch}

\begin{proposition}[Parameter-dependent reset triggers break continuity]
\label{prop:gsv-ess-branch}
Fix $\tau\in(0,1)$ and let the program apply the reset at time $t$ if and
only if $\mathrm{ESS}_t(\theta)=\big(\sum_i\alpha_{t,i}^{2}\big)^{-1}<\tau N$.
Suppose at $\theta^{*}$ some time $t^{*}$ satisfies
$\mathrm{ESS}_{t^{*}}(\theta^{*})=\tau N$ with
$\nabla_\theta\,\mathrm{ESS}_{t^{*}}(\theta^{*})\neq0$, and that the two
branch programs (reset applied at $t^{*}$ versus skipped) produce different
values of some later increment at $\theta^{*}$.  Then
$\theta\mapsto\widehat L^{N}(\theta)$ is discontinuous at $\theta^{*}$, and
no same-scalar derivative claim holds across the crossing.
\end{proposition}

\begin{proof}
By the implicit function theorem the crossing set is locally a hypersurface
separating two open regions; on one side the executed program is the
reset-at-$t^{*}$ branch, on the other the skip branch.  Each branch value is
continuous on its closed side, and their limits at $\theta^{*}$ are the two
branch values at $\theta^{*}$, which differ by hypothesis (the carried
clouds differ from $t^{*}$ on, and by assumption some later increment
differs).  A function with unequal one-sided limits is discontinuous.  The
final statement follows because a total derivative of $\widehat L^{N}$ at
$\theta^{*}$ requires continuity there.  Consequently $\theta$-independent
reset schedules are the safe policy for the repository's fixed-branch score
contract.  Parameter-dependent branches are not mathematically impossible in
general, but would require separately verified value and derivative matching at
every boundary (or a smooth gate).
\end{proof}

\section{Claim ledger and non-claims}
\label{sec:gsv-claims}

\emph{Proved here:}
Propositions~\ref{prop:gsv-tangent-recursion}--\ref{prop:gsv-gn-gain}
(recursion, increment tangent, stage-gain bounds);
Proposition~\ref{prop:gsv-cubature-design-moments} (design defect
identities); Propositions~\ref{prop:gsv-whitened-design-identities},
\ref{prop:gsv-iid-gaussian-moments}, \ref{prop:gsv-delta-method},
\ref{prop:gsv-genut-positivity} and
Corollary~\ref{cor:gsv-design-tangent} (design repair);
Propositions~\ref{prop:gsv-leapfrog-volume}--\ref{prop:gsv-mh-invariance}
and Corollary~\ref{cor:gsv-surrogate-force} (surrogate-force HMC);
Propositions~\ref{prop:gsv-fisher-identity}--\ref{prop:gsv-backward-recursion}
(Fisher identity and backward recursion);
Proposition~\ref{prop:gsv-ess-branch} (branch warning).

\emph{Cited, not re-verified locally:} the $O(T)$ versus $O(T^{2})$
variance-growth rates of Remark~\ref{rem:gsv-pds-boundary}; the primary
Poyiadjis--Doucet--Singh PDF is not yet available as a valid local copy.

\emph{Hypotheses, explicitly not theorems:} a positive finite-time directional
growth summary on the Austria diagonal-only route
(Remark~\ref{rem:gsv-lyapunov}); the causal attribution of the observed
score variance to the design defect
(Remark~\ref{rem:gsv-defect-consequence}); the $O_p(N^{-1})$
whitening-versus-studentization gap (Remark~\ref{rem:gsv-whitening-gap}).

\emph{Non-claims:} nothing here establishes an exact-posterior score,
unbiasedness of $\widehat L^{N}$, variance superiority of any repair,
acceptance rates or mixing of surrogate-force HMC, or any default, HMC, or
leaderboard promotion.  Every proposed run requires its own experiment plan
with an evidence contract, fresh tuning scope, and untouched claim seeds.

\section{Code and artifact anchors (non-mathematical)}
\label{sec:gsv-anchors}

\begin{center}
\small
\begin{tabular}{@{}p{0.42\textwidth}p{0.52\textwidth}@{}}
\toprule
Object in this note & Repository anchor \\
\midrule
Tangent recursion \eqref{eq:gsv-recursion} &
\texttt{cubature\_genut\_filter.py}, \texttt{time\_body} (466--698) \\
Increment tangent \eqref{eq:gsv-increment-tangent} &
\texttt{cubature\_genut\_filter.py} (533--544) \\
Affine tangent \eqref{eq:gsv-affine-tangent} &
\texttt{ledh\_contract\_e\_reset\_tf.py} (319--328) \\
GN coefficient tangent \eqref{eq:gsv-gn-tangent} &
\texttt{higher\_moment\_contract\_e.py} (421--458) \\
Cubature design of
Prop.~\ref{prop:gsv-cubature-design-moments} &
\texttt{cubature\_genut\_candidate.py}, \texttt{cubature\_design} (129--144) \\
GenUT positivity boundary &
\texttt{cubature\_genut\_candidate.py},
\texttt{genut\_design}/\texttt{replicate\_positive\_genut} \\
Chapter score proposition &
\texttt{docs/chapters/ch32c\_entropic\_ot\_sinkhorn.tex} (2342--2371) \\
Austria empirical anchors &
\texttt{docs/plans/bayesfilter-austria-sir-pairwise-moment-genut-score-trial-result-2026-07-30.md} \\
2026-07-30 code audit &
\texttt{docs/plans/bayesfilter-genut-score-computation-audit-result-2026-07-30.md} \\
\bottomrule
\end{tabular}
\end{center}

\end{document}
